\documentclass{article}

\usepackage[dvipsnames]{xcolor}

\usepackage{tikz-cd}

\usepackage[utf8]{inputenc}
\usepackage{amsmath}
\usepackage{amsfonts}
\usepackage{amssymb}
\usepackage[margin=.75in, total={8.25in, 10.75in}, heightrounded]{geometry}

\usepackage{graphicx}
\graphicspath{{graphics/}}

\usepackage{scalerel}

\usepackage{stmaryrd}

\usepackage{MnSymbol}

\usepackage{mdframed}

\usepackage{titlesec}

\usepackage{enumitem}

\usepackage{stackengine}

\setenumerate[0]{label=\alph*)}

\usepackage{blkarray}

\usepackage{mathtools}

\usepackage{hyperref}

\hypersetup
{
	colorlinks = true,
	urlcolor = OliveGreen
}

\titleformat{\section}
{\normalfont \Large \bfseries \centering}{\Roman{section} --- }{0pt}{}



\newcommand{\hooklongrightarrow}{\lhook\joinrel\longrightarrow}
\newcommand{\hooklongleftarrow}{\longleftarrow\joinrel\rhook}


\renewcommand{\baselinestretch}{1.33}



\definecolor{DefGreen}{rgb}{0,0.5,0}
\definecolor{TheoremOrange}{rgb}{0.88,0.6,0.08}
\definecolor{LemmaYellow}{rgb}{1,1,0}
\definecolor{CorollaryBlue}{rgb}{0,0.3,1}
\definecolor{ProofPurple}{rgb}{0.58,0,1}
\definecolor{AxiomRed}{rgb}{1,0,0}
\definecolor{CommentBlue}{rgb}{0.46,0.67,1}



\usepackage{amsthm}

\newtheoremstyle{colontheorem}
	{0in}                    	% Space above
	{.15in}                   	% Space below
	{\normalfont}      		    % Body font
	{}                          % Indent amount
	{\bfseries}                 % Theorem head font
	{:}                         % Punctuation after theorem head
	{.5em}                      % Space after theorem head
	{}							% Theorem head spec (can be left empty, meaning ‘normal’)
	
\theoremstyle{colontheorem}

\newtheorem{theorem}{Theorem}[section]
\newtheorem{proposition}[theorem]{Proposition}
\newtheorem{definition}[theorem]{Definition}
\newtheorem{axiom}[theorem]{Axiom}

\newtheorem{lemma}{Lemma}[theorem]
\newtheorem{corollary}{Corollary}[theorem]

\newtheorem{exercise}{Exercise}




\newcommand{\fadeline}
{
	\noindent\begin{tikzpicture}[baseline]
		\path[left color=white,right color=white,middle color=black]
		(0,0) rectangle (\textwidth,.5pt);%
	\end{tikzpicture}
}

\newcommand{\Span}{\textnormal{span}}
\newcommand{\Null}{\textnormal{null }}
\newcommand{\Range}{\textnormal{range }}
\newcommand{\T}{^\textnormal{T}}

\newcommand{\Sub}{\textnormal{sub }}

\newcommand{\re}{\textnormal{Re }}
\newcommand{\im}{\textnormal{Im }}
\newcommand{\Arg}{\textnormal{Arg }}
\newcommand{\Log}{\textnormal{Log }}
\newcommand{\Res}{\textnormal{Res}}
\newcommand{\pv}{\textnormal{p.v.}}

\newcommand{\e}{\varepsilon}

\newcommand{\Aut}{\textnormal{Aut }}
\newcommand{\Inn}{\textnormal{Inn }}
\newcommand{\orb}{\textnormal{orb }}
\newcommand{\Char}{\textnormal{char }}
\newcommand{\Spec}{\textnormal{Spec }}
\newcommand{\Frac}{\textnormal{Frac }}
\newcommand{\Gal}{\textnormal{Gal}}

\newcommand{\transverse}{\stackinset{c}{}{c}{.75pt}{\rotatebox[origin=c]{180}{$\perp$}}{$\scaleto{\cap}{1ex}$}}





\newenvironment{Theorem}
{
	\begin{mdframed}[backgroundcolor=TheoremOrange!10]
	\begin{theorem}
}
{
	\end{theorem}
	\end{mdframed}
	
	\vspace{.15in}
}

\newenvironment{Proposition}
{
	\begin{mdframed}[backgroundcolor=TheoremOrange!10]
	\begin{proposition}
}
{
	\end{proposition}
	\end{mdframed}
	
	\vspace{.15in}
}

\newenvironment{Def}
{
	\begin{mdframed}[backgroundcolor=DefGreen!10]
	\begin{definition}
}
{
	\end{definition}
	\end{mdframed}
	
	\vspace{.15in}
}

\newenvironment{Axiom}
{
	\begin{mdframed}[backgroundcolor=AxiomRed!10]
	\begin{axiom}
}
{
	\end{axiom}
	\end{mdframed}
	
	\vspace{.15in}
}

\newenvironment{Lemma}
{
	\begin{mdframed}[backgroundcolor=LemmaYellow!10]
	\begin{lemma}
}
{
	\end{lemma}
	\end{mdframed}
	
	\vspace{.03in}
}

\newenvironment{Corollary}
{
	\vspace{-.035in}
	
	\begin{mdframed}[backgroundcolor=CorollaryBlue!10]
	\begin{corollary}
}
{
	\end{corollary}
	\end{mdframed}
	
	\vspace{.09in}
}

\newenvironment{Comment}
{
	\begin{mdframed}[backgroundcolor=CommentBlue!10]
	\textbf{Comment:}%
}
{
	\end{mdframed}
	
	\vspace{.15in}
}

\newenvironment{Proof}
{
	\begin{mdframed}[backgroundcolor=ProofPurple!10]
	\textbf{Proof:}%
}
{
	\end{mdframed}
	
	\vspace{.085in}
}

\newenvironment{Example}
{
	\begin{mdframed}
	\textbf{Example:}%
}
{
	\end{mdframed}
	
	\vspace{.15in}
}

\newenvironment{Exercise}
{
	\begin{mdframed}
	\begin{exercise}
}
{
	\end{exercise}
	\end{mdframed}
	
	\vspace{.15in}
}




\setlength{\parindent}{0pt}


\begin{document}

\vspace*{.5in}

\begin{center}
	\Huge Applied Notes\\
	
	\vspace{.25in}
	
	\normalsize February 12--14, 2020
\end{center}

\vspace{.25in}



\begin{Comment}
	%
	Let $S(t)$ be the CMTP on $\mathbb{Z}$ that jumps up and down by one, both at rate $\frac{1}{2}$. What does $S$ look like in the long term? It always has mean $0$, and so the variance is just $\mathbb{E}[S(t)^2] = \mathbb{E}[(U(t) - D(t))^2]$, where $U(t)$ is the number of jumps up before time $t$ and $D(t)$ is the number of jumps down. Since both $U(t)$ and $S(t)$ are Poisson with rate $\frac{t}{2}$, the variance of $S(t)$ is just $t$. Thus the standard deviation is $\sqrt{t}$. Fixing $N$, let $B_N(t) = \frac{1}{\sqrt{N}} S(Nt)$. Then $\mathbb{E}[B_N(t)] = 0$ and $\text{var}(B_N(t)) = t$. This is called Brownian Scaling. The question is, does $B_N(t)$ converge in distribution to some $B(t)$ as $N \longrightarrow \infty$?
	
\end{Comment}



\begin{Lemma}
	
	If $X_N$ and $X$ are Markov processes such that for all $s$ with $s < t$, all $x_0$, and all sets $A$,
	
	$$
		\mathbb{P}(X_N(t) \in A \mid X_N(s) = x_0) \longrightarrow \mathbb{P}(X(t) \in A \mid X(s) = x_0),
	$$
	
	then $X_N \longrightarrow X$ in distribution.
	
\end{Lemma}



\begin{Def}
	
	The $B_N$ converge to some $B$ by the lemma, and it is called \textbf{Brownian Motion}.
	
\end{Def}



\begin{Theorem}
	
	If $P(t, x) = \mathbb{P}(B(t) = x \mid B(0) = 0)$, then $\frac{d}{dt} [P(t, x)] = \frac{1}{2} \frac{d^2}{dx^2} [P(t, x)]$, $P(0, x) = \delta_0(x)$, and $P(t, x) = \dfrac{1}{\sqrt{2 \pi t}} e^{-x^2 / 2t}$ --- that is, $B(t) \sim N(0, t)$.
	
\end{Theorem}



\begin{Lemma}
	
	If $W \sim N(0, t)$, then $\mathbb{E}[e^{-uW}] = e^{-u^2 t / 2}$. (This is the Laplace Transform, which uniquely determines probability distributions.)
	
\end{Lemma}



\begin{Proof}
	%
	$$
		\mathbb{E}[e^{-uW}] = \int_{-\infty}^\infty e^{-ux} \cdot \frac{1}{\sqrt{2 \pi t}} e^{-x^2 / 2t}\ dx = e^{-u^2 t / 2} \int_{-\infty}^\infty \frac{1}{\sqrt{2 \pi t}} e^{-(x + ut)^2 / 2t} dx
	$$
	
	by completing the square. But that integral is just a shifted normal distribution, so it evaluates to $1$. Thus $\mathbb{E}[e^{-uW}] = e^{-u^2 t / 2}$.
	
\end{Proof}



\begin{Comment}
	%
	We can also construct $B$ another way. Let $\Lambda$ be a PPP on $[0, \infty) \times \mathbb{R}$, where $[0, \infty)$ represents time and $\mathbb{R}$ represents jump size, with mean intensity $dt \cdot f(x) \cdot dx$, where $f$ is a probability density (that is, $\int_{-\infty}^\infty f(x)\ dx = 1$). Now if $X$ has probability density $f$, then $\mathbb{E}[X] = 0 = \int_{-\infty}^\infty x f(x)\ dx$ and $\mathbb{E}[X^2] = \sigma^2 = \int_{-\infty}^\infty x^2 f(x)\ dx$. Let $S(t)$ be the sum of the jumps out to time $t$. Then
	
	$$
		S(t) = \sum_{\substack{(t_i, x_i) \in \Lambda \\ t_i \leq t}} x_i = \int_0^t \int_{-\infty}^\infty x \Lambda(dt, dx).
	$$
	
	Now let $B_N(t) = \frac{1}{\sqrt{N}} S(Nt) = \frac{1}{\sqrt{N}}\int_0^t \int_{-\infty}^\infty x \Lambda(dt, dx)$. Then
	
	\begin{align*}
		\mathbb{E}[e^{-u B_N(t)}] &= \exp \left( \int_0^{tN} \int_{-\infty}^\infty (e^{\frac{ux}{\sqrt{N}}} - 1) f(x)\ ds\ dx \right)\\
		&= \exp \left( \int_0^{tN} \mathbb{E} \left[ e^{\frac{uX}{\sqrt{N}}} - 1 \right]\ ds \right)\\
		&= \exp \left( \int_0^{tN} \mathbb{E} \left[ \frac{uX}{\sqrt{N}} + \frac{u^2 X^2}{2 \sqrt{N}} + O\left( \frac{1}{N^{3/2}} \right) \right]\ ds \right)\\
		&= \exp \left( \int_0^{tN} 0 + \frac{u^2 \sigma^2}{2 \sqrt{N}} + O\left( \frac{1}{N^{3/2}} \right)\ ds \right)\\
		&= e^{- \frac{1}{2} \sigma^2 u^2 t} + O\left( \frac{1}{\sqrt{N}} \right) \sim N(0, \sigma^2 t).
	\end{align*}
	
	Thus $B_N \longrightarrow B(\sigma^2, t)$ in distribution, since $B_N$ is a Markov process and has iid increments --- in particular, $B_N(t) - B_N(s) = B_N(t - s)$ in distribution.
	
\end{Comment}



\begin{Theorem}
	
	1. $B$ is a standard Brownian motion if and only if $B$ is a Markov process and $B(t) - B(s) \sim N(0, t - s)$ for all $0 \leq s < t$.
	
	2. The map $t \longmapsto B(t)$ is continuous but nowhere differentiable almost surely.
	
	3. If $W(t) = \dfrac{1}{\sqrt{c}} B(ct)$ for some $c$, then $W(t)$ and $B(t)$ are equal in distribution.
	 	
\end{Theorem}



\end{document}
